%%%%%%%%%%%%%%%%%%%%%%%%%%%%%%%%%%%%%%%%%%%%%%%%%%%%%%%%%%%%%%%%%%%%%%
% Lớp 8 | Chương I. Đa thức | Đơn thức
% Dạng 8D1?1-1: Nhận biết đơn thức, đơn thức thu gọn, bậc của đơn thức
%%%%%%%%%%%%%%%%%%%%%%%%%%%%%%%%%%%%%%%%%%%%%%%%%%%%%%%%%%%%%%%%%%%%%%


% nguồn: De_AI_Toan8_cuoi_hoc_ki_i_1 (nhập 21/11/2026)
\begin{bt}%[8D1N1-1]%[ma:4ea5c559]
	Trong các biểu thức sau, biểu thức nào là đơn thức? Chỉ ra hệ số, phần biến và bậc của các đơn thức đó:
$3x^2y$; $-5$; $x + 2y$; $\dfrac{1}{2}xy^3z$; $\dfrac{2}{x}$.
	\loigiai{
		Các biểu thức là đơn thức gồm: $3x^2y$, $-5$, $\dfrac{1}{2}xy^3z$ (các biểu thức $x + 2y$ và $\dfrac{2}{x}$ không phải là đơn thức). \\
- Đơn thức $3x^2y$ có hệ số $3$, phần biến $x^2y$, bậc $2 + 1 = 3$. \\
- Đơn thức $-5$ có hệ số $-5$, không có phần biến, bậc $0$. \\
- Đơn thức $\dfrac{1}{2}xy^3z$ có hệ số $\dfrac{1}{2}$, phần biến $xy^3z$, bậc $1 + 3 + 1 = 5$.
	}
\end{bt}

% nguồn: De_AI_Toan8_cuoi_hoc_ki_i_1 (nhập 21/11/2026)
\begin{bt}%[8D1N1-1]%[ma:8bff8ce5]
	Thu gọn các đơn thức sau rồi xác định hệ số, phần biến và bậc của mỗi đơn thức:
a) $A = 2x^2y \cdot (-3xy^3)$;
b) $B = \dfrac{3}{4}xy^2 \cdot \left(-\dfrac{8}{9}x^3y\right)$.
	\loigiai{
		a) $A = 2 \cdot (-3) \cdot (x^2 \cdot x) \cdot (y \cdot y^3) = -6x^3y^4$. \\
Hệ số là $-6$, phần biến là $x^3y^4$, bậc là $3 + 4 = 7$. \\
b) $B = \dfrac{3}{4} \cdot \left(-\dfrac{8}{9}\right) \cdot (x \cdot x^3) \cdot (y^2 \cdot y) = -\dfrac{2}{3}x^4y^3$. \\
Hệ số là $-\dfrac{2}{3}$, phần biến là $x^4y^3$, bậc là $4 + 3 = 7$.
	}
\end{bt}

% nguồn: De_AI_Toan8_cuoi_hoc_ki_i_1 (nhập 21/11/2026)
\begin{bt}%[8D1N1-1]%[ma:0d80cc84]
	Thu gọn các đơn thức sau rồi tìm bậc của mỗi đơn thức:
a) $M = (-2xy^2)^2 \cdot 3x^2y$;
b) $N = \left(-\dfrac{1}{2}x^2y\right)^3 \cdot 8xy^2$.
	\loigiai{
		a) $M = 4x^2y^4 \cdot 3x^2y = 12x^4y^5$. Bậc của $M$ là $4 + 5 = 9$. \\
b) $N = \left(-\dfrac{1}{8}x^6y^3\right) \cdot 8xy^2 = -x^7y^5$. Bậc của $N$ là $7 + 5 = 12$.
	}
\end{bt}

% nguồn: De_AI_Toan8_cuoi_hoc_ki_i_1 (nhập 21/11/2026)
\begin{bt}%[8D1N1-1]%[ma:5c5124c4]
	Thu gọn các đơn thức ba biến sau và cho biết hệ số, bậc của mỗi đơn thức:
a) $P = 5xyz \cdot (-4x^2yz^3)$;
b) $Q = \dfrac{2}{5}x^2yz \cdot \dfrac{5}{6}xy^3z^2$.
	\loigiai{
		a) $P = 5 \cdot (-4) \cdot x^3y^2z^4 = -20x^3y^2z^4$. Hệ số là $-20$, bậc là $3 + 2 + 4 = 9$. \\
b) $Q = \left(\dfrac{2}{5} \cdot \dfrac{5}{6}\right)x^3y^4z^3 = \dfrac{1}{3}x^3y^4z^3$. Hệ số là $\dfrac{1}{3}$, bậc là $3 + 4 + 3 = 10$.
	}
\end{bt}

% nguồn: De_AI_Toan8_cuoi_hoc_ki_i_1 (nhập 21/11/2026)
\begin{bt}%[8D1N1-1]%[ma:fc88b6b0]
	Thu gọn đa thức $A = -2xy + 4xy^2 + xy + xy^2$ rồi xác định các hạng tử và bậc của đa thức $A$.
	\loigiai{
		Ta nhóm các hạng tử đồng dạng: $A = (-2xy + xy) + (4xy^2 + xy^2) = -xy + 5xy^2$. \\
Đa thức $A$ có hai hạng tử là $-xy$ (bậc $2$) và $5xy^2$ (bậc $3$). \\
Vậy bậc của đa thức $A$ là $3$.
	}
\end{bt}

% nguồn: De_AI_Toan8_cuoi_hoc_ki_i_1 (nhập 21/11/2026)
\begin{bt}%[8D1N1-1]%[ma:187f5749]
	Thu gọn đa thức $M = 3x^2y - 2xy^2 + x^3 - x^2y + xy^2 - x^3$ rồi tìm bậc của đa thức $M$.
	\loigiai{
		Ta có $M = (3x^2y - x^2y) + (-2xy^2 + xy^2) + (x^3 - x^3) = 2x^2y - xy^2$. \\
Cả hai hạng tử $2x^2y$ và $-xy^2$ đều có bậc $3$, nên bậc của đa thức $M$ là $3$.
	}
\end{bt}

% nguồn: De_AI_Toan8_cuoi_hoc_ki_i_1 (nhập 21/11/2026)
\begin{bt}%[8D1N1-1]%[ma:42c426d0]
	Thu gọn đa thức $D = 3x^2y^2 - 5xy + x^2y^2 + 2xy - 4x^2y^2 + 1$ rồi xác định bậc của $D$.
	\loigiai{
		Ta có $D = (3x^2y^2 + x^2y^2 - 4x^2y^2) + (-5xy + 2xy) + 1 = -3xy + 1$. \\
Hạng tử $-3xy$ có bậc $2$, hạng tử $1$ có bậc $0$, nên bậc của đa thức $D$ là $2$.
	}
\end{bt}

% nguồn: De_AI_Toan8_cuoi_hoc_ki_i_1 (nhập 21/11/2026)
\begin{bt}%[8D1N1-1]%[ma:d9c12990]
	Thu gọn đa thức $B = xy^2z + 2xy^2z - xyz - 3xy^2z + xy^2z$ rồi tìm bậc của $B$.
	\loigiai{
		Ta có $B = (xy^2z + 2xy^2z - 3xy^2z + xy^2z) - xyz = xy^2z - xyz$. \\
Hạng tử $xy^2z$ có bậc $1 + 2 + 1 = 4$, hạng tử $-xyz$ có bậc $3$, nên bậc của đa thức $B$ là $4$.
	}
\end{bt}

% nguồn: De_AI_Toan8_cuoi_hoc_ki_i_1 (nhập 21/11/2026)
\begin{bt}%[8D1N1-1]%[ma:eb010976]
	Thu gọn đa thức $C = \dfrac{1}{2}xyx - 2xy + x - x^2y + 3xy - 4x$ rồi tìm bậc của $C$.
	\loigiai{
		Ta viết $\dfrac{1}{2}xyx = \dfrac{1}{2}x^2y$. \\
Suy ra $C = \left(\dfrac{1}{2}x^2y - x^2y\right) + (-2xy + 3xy) + (x - 4x) = -\dfrac{1}{2}x^2y + xy - 3x$. \\
Hạng tử có bậc cao nhất là $-\dfrac{1}{2}x^2y$ (bậc $3$), vậy bậc của đa thức $C$ là $3$.
	}
\end{bt}
